\section{Quality and Validation}
\label{sec:quality}

MOCHO was built with the Ontology Development Kit (ODK).
The ontology requires OWL~2~DL expressivity: it makes use of universal restrictions
(\texttt{owl:allValuesFrom}), cardinality constraints, complement and disjointness axioms,
and property chains, which collectively place it outside the OWL~2~EL profile.

OWL~2~DL consistency is verified by HermiT~\cite{motik2009hypertableau} as part of the
ODK build; materialized inferences are exported as RDF triples and loaded into QLever
alongside the DDB corpus, so no runtime reasoning is performed by the triplestore.
The ontology and its documentation are published through the ODK and
Widoco~\cite{garijo2017widoco} pipeline, producing both the machine-readable artifact and
human-readable HTML documentation from a single build step.
The EDM-to-MOCHO property mappings are validated by a suite of unit tests that assert
correct predicate assignment per \texttt{(target\_class,~ddb-edm\_prop)} pair against
gold-standard triples derived from the Goethe-Faust corpus.

\subsection{Completeness}
\label{sec:quality:completeness}

GEMEA is a snapshot exported between late March and early April~2026, covering
97.7\,\% of the digitized DDB corpus at time of writing (Table~\ref{tab:scale}).
The remaining gap reflects objects added to the DDB after the snapshot window
and a small number of transient API search failures during the export run.

\subsection{Scope of Validation}
\label{sec:quality:scope}

Comprehensive metadata quality evaluation of a cultural heritage
knowledge graph at this scale presents a structural bootstrapping problem: rigorous evaluation requires
the dataset to be published and accessible, yet conducting and reporting such an evaluation
prior to publication is itself impractical at this scale.
GEMEA is released to enable this evaluation; the SPARQL endpoint, Linked Data browser,
and downloadable dump provide the infrastructure for systematic quality studies at corpus
scale.

\subsection{Known Limitations}
\label{sec:quality:limitations}

\begin{itemize}
  \item \textbf{Snapshot currency.} The dataset reflects the DDB as of late March to
        early April~2026. Records added, updated, or deleted after the export window
        are not reflected; no incremental update mechanism is currently implemented.

  \item \textbf{Partial GND linking.} Werk linking covers literary and musical works but achieves low recall due to
        inconsistent encoding of \texttt{dc:title} values across contributing institutions:
        corpus analysis of the Goethe-Faust reference corpus shows 71\,\% of titles
        carry no ISBD punctuation~\cite{ifla2012isbd}, with the dominant signal being
        the subtitle separator `\,:\,' (18\,\%) rather than the statement-of-responsibility
        separator `\,/\,' (2.1\,\%) used for work-title extraction; corpus-wide figures
        for 9.2M German-language DDB titles show 92.4\,\% without any ISBD marker.
        Future work will apply ISBD punctuation heuristics and a planned NER model
        fine-tuned on a DDB subset to improve coverage; authority linking for images
        of artworks and architecture via Wikidata,
        Bildindex\footnote{\url{https://www.bildindex.de}}, and artresearch.net's
        PHAROS infrastructure
        is also planned.

  \item \textbf{Alignment coverage at scale.} The EDM-to-MOCHO mapping was developed
        and validated against the Goethe-Faust corpus (115,432 library records).
        Sector-specific edge cases, particularly in the Library sector at 18M records,
        may produce fallback or underspecified type assignments.

  \item \textbf{Rights metadata.} Each object retains the \texttt{dcterms:rights}
        value as declared by the contributing institution; rights statements are not
        normalized, verified, or mapped to a controlled vocabulary.

  \item \textbf{Entity resolution for bare DDB identifiers.} Agents, places, concepts,
        and time spans that lack a GND or other authority URI are assigned
        DDB-internal bare identifiers, which are record-scoped and not shared across
        contributing institutions. The same real-world entity may therefore appear as
        multiple distinct nodes in the graph. Authority-based entity resolution for
        bare-identifier entities is left as future work.

  \item \textbf{Geographic and temporal sparsity.} Corpus analysis of the Goethe-Faust
        reference corpus shows that only 45\,\% of objects carry any
        date value (\texttt{dc:date}, \texttt{dcterms:date}, or \texttt{dcterms:created})
        and only 7.1\,\% carry a place value (\texttt{dc:coverage} or
        \texttt{dcterms:spatial}).
\end{itemize}
